\subsubsection{Pollard-Rho}
Pollard's Rho 算法，一种用于大整数质因数分解的随机算法。
\begin{lstlisting}
ll qp(ll a,ll b,ll p){a%=p;
	ll ans=1;while(b){
		if(b&1)ans=(__int128)ans*a%p;
		a=(__int128)a*a%p;b>>=1;
	}return ans;
}ll n,m;
ll tst[7]={2,325,9375,28178,450775,9780504,1795265022};
/*
unsigned : 2 7 61
ull : 2 325 9375 28178 450775 9780504 1795265022
*/
bool miller(ll x){if(x<2)return 0;
    if(x==2||x==3||x==5||x==7||x==11||x==13)return 1;
	ll u=x-1,t=0;while(u%2==0)u/=2,t++;
	f(T,0,6){
		ll a=tst[T];
		if(x==a)return 1;
		ll b=qp(a,u,x),bb;//if(b==1||b==x-1)continue;
		f(i,1,t){
			bb=(__int128)b*b%x;
			if(bb==1&&b!=1&&b!=x-1)return 0;b=bb;
		}if(b!=1)return 0;
	}return 1;
}ll pollard(ll p){
	ll c=(ll)rand()%p+1;
	ll s=0,t=0,mul=1,k=1,i=0;
	for(k=1;;k<<=1,s=t,mul=1){
		for(i=1;i<=k;i++){
			t=((__int128)t*t+c)%p;
			mul=(__int128)mul*abs(s-t)%p;
			if(i%127==0){
				ll g=__gcd(mul,p);
				if(g>1)return g;
			}
		}ll g=__gcd(mul,p);
		if(g>1)return g;
	}
}ll mxp;
void work(ll p){
	if(p<2||p<=mxp)return;
	if(miller(p)){mxp=max(mxp,p);return;}
	ll k=p;while(k==p)k=pollard(p);
	while(p%k==0)p/=k;
	work(p),work(k);
}
int main(){srand(time(0));
	wt{	ll p=d;mxp=0;
		work(p);if(mxp==p)puts("Prime");
		else printf("%lld\n",mxp);
	}
	return 0;
}
\end{lstlisting}

